% ============================================================
%  BEAMER — CLASE DE APOYO: ARITMÉTICA
%  Breve historia de los números · Conjuntos numéricos ·
%  Potenciación · Operaciones con fracciones · Fracciones
%  generatrices · Fracción irreducible
%  Autor: Ing. Gabriel Astudillo
%  Compilar: latexmk -xelatex beamer_apoyo_aritmetica.tex
%  (o desde la raíz: ./compile_all.sh beamer_apoyo_aritmetica)
% ============================================================
\documentclass[aspectratio=169]{beamer}
\usepackage{mathwizards-beamer}
\mwbeamer{Clase de Apoyo}{Aritmética: Números, Conjuntos Numéricos y Fracciones}{3er Año}

\begin{document}

\begin{frame}[plain]
  \titlepage
\end{frame}

% ════════════════════════════════════════════════════════════
\section{Lo que Veremos Hoy}
% ════════════════════════════════════════════════════════════

\begin{frame}{Hoja de Ruta de la Clase}
  \begin{columns}[T]
    \column{0.48\textwidth}
    \begin{itemize}
      \item Breve historia de los números
      \item Conjuntos numéricos: $\mathbb{N}$, $\mathbb{Z}$ y $\mathbb{Q}$
      \item Potenciación en $\mathbb{Q}$
      \item Suma y resta de fracciones
      \item Producto y división de fracciones
    \end{itemize}
    \column{0.48\textwidth}
    \begin{itemize}
      \item Fracciones generatrices
      \item Fracción irreducible
      \item Números reales y aproximaciones
      \item Práctica: operaciones combinadas
      \item \textbf{Ejercicios de repaso general}
    \end{itemize}
  \end{columns}
  \vspace{4pt}
  \begin{mwpasos}[Cómo trabajaremos]
    Cada tema termina con \textbf{ejercicios propuestos} y con su
    \textbf{solución explicada paso a paso}: resuélvelos primero a mano
    sobre el PDF y luego compara tu desarrollo.
  \end{mwpasos}
\end{frame}

% ════════════════════════════════════════════════════════════
\section{Breve Historia de los Números}
% ════════════════════════════════════════════════════════════

\begin{frame}{Marcas de Cuentas y Numeraciones Antiguas}
  \begin{block}{Marcas de cuentas}
    El método más antiguo para contar: se agrupan marcas y se cuenta
    cuántas hay.
    \[
      \underbrace{|\,|\,|\,|}_{4}\;\;\underbrace{|\,|\,|\,|}_{4} \;=\; 8
    \]
  \end{block}
  \vspace{4pt}
  \begin{columns}[T]
    \column{0.48\textwidth}
    \begin{mwextra}[Números babilonios]
      \textbf{Sistema posicional}: el valor de un carácter depende de
      su \textbf{posición}. Base sexagesimal (60).
    \end{mwextra}
    \column{0.48\textwidth}
    \begin{mwextra}[Números egipcios]
      Sistema \textbf{aditivo}: los valores de cada símbolo se
      \textbf{suman}, sin importar la posición.
    \end{mwextra}
  \end{columns}
\end{frame}

\begin{frame}{De Roma a la India}
  \begin{columns}[T]
    \column{0.48\textwidth}
    \begin{block}{Números romanos}
      Notación \textbf{sustractiva}: un numeral menor a la izquierda
      se \textbf{resta} del mayor.
      \[
        \text{IIII} \to \text{IV}
        \qquad
        \text{VIIII} \to \text{IX}
      \]
      \textbf{Limitación}: poco aptos para el cálculo aritmético.
    \end{block}
    \column{0.48\textwidth}
    \begin{block}{Números arábigo-hindúes}
      Herencia de Babilonia e India (siglo XI): sistema
      \textbf{posicional decimal} con el \textbf{cero}.
      \[
        1405 = 1\cdot10^{3}+4\cdot10^{2}+0\cdot10+5
      \]
      Es el sistema que usamos \textbf{hoy}.
    \end{block}
  \end{columns}
  \vspace{4pt}
  \begin{mwextra}[Legado moderno del sistema sexagesimal]
    Seguimos usando la base 60: \textbf{60 minutos} por hora,
    \textbf{60 segundos} por minuto, \textbf{360°} en una vuelta
    completa.
  \end{mwextra}
\end{frame}

\begin{frame}{Ejercicios Propuestos — Historia de los Números \textnormal{(1/2)}}
  \begin{mwejercicio}[Cuestionario 1--4 \quad \facil]
    \begin{enumerate}
      \item ¿Qué importante hallazgo arqueológico está considerado
            como una de las evidencias más antiguas del uso de marcas
            numéricas?
      \item ¿Qué característica distingue fundamentalmente al sistema
            numérico inventado por Babilonia?
      \item ¿En qué legado moderno se observa aún la influencia del
            sistema sexagesimal?
      \item ¿Cuál fue la principal limitación de la numeración romana
            para el cálculo aritmético?
    \end{enumerate}
  \end{mwejercicio}
\end{frame}

\begin{frame}{Ejercicios Propuestos — Historia de los Números \textnormal{(2/2)}}
  \begin{mwejercicio}[Cuestionario 5--8 \quad \facil]
    \begin{enumerate}
      \setcounter{enumi}{4}
      \item ¿Cuál fue la causa histórica por la cual fue tan versátil
            la adaptación del sistema decimal?
      \item ¿Qué crédito tuvo India durante el siglo VII en el
            desarrollo de la numeración y los signos operativos?
      \item ¿Cómo interpreta Brahmagupta la noción intuitiva de los
            números negativos?
      \item ¿Qué importante centro cultural del imperio selyúcida
            sirvió de puente entre India y Europa?
    \end{enumerate}
  \end{mwejercicio}
\end{frame}

% ════════════════════════════════════════════════════════════
\section{Conjuntos Numéricos: \texorpdfstring{$\mathbb{N}$}{N},
  \texorpdfstring{$\mathbb{Z}$}{Z} y \texorpdfstring{$\mathbb{Q}$}{Q}}
% ════════════════════════════════════════════════════════════

\begin{frame}{Conjunto de los Números Naturales}
  \begin{block}{Definición}
    \[
      \mathbb{N} = \{0,\, 1,\, 2,\, 3,\, \dots\}
    \]
    Son números \textbf{enteros}, \textbf{discretos} e
    \textbf{infinitos}.
  \end{block}
  \vspace{4pt}
  \begin{mwextra}[Propiedades de la suma en $\mathbb{N}$]
    \begin{itemize}
      \item \textbf{Conmutativa:} $a + b = b + a$
      \item \textbf{Asociativa:} $(a + b) + c = a + (b + c)$
      \item \textbf{Elemento neutro:} $a + 0 = a$
    \end{itemize}
  \end{mwextra}
\end{frame}

\begin{frame}{Conjunto de los Números Enteros}
  \begin{block}{Definición}
    \[
      \mathbb{Z} = \{\dots,\, -3,\, -2,\, -1,\, 0,\, 1,\, 2,\, 3,\, \dots\}
    \qquad \mathbb{N} \subset \mathbb{Z}
    \]
    \textbf{Particiones:}\quad
    $\mathbb{Z}^{+} = \mathbb{N}$ \quad$\cdot$\quad
    $\mathbb{Z}^{-} = -\mathbb{N}$ \quad$\cdot$\quad
    $\mathbb{Z}^{*} = \mathbb{Z} \setminus \{0\}$
  \end{block}
  \begin{columns}[T]
    \column{0.48\textwidth}
    \begin{mwextra}[Suma en $\mathbb{Z}$]
      Además de las propiedades de $\mathbb{N}$, existe el
      \textbf{elemento inverso aditivo}:
      $a + (-a) = a - a = 0$.
    \end{mwextra}
    \column{0.48\textwidth}
    \begin{mwextra}[Producto en $\mathbb{Z}$]
      \begin{itemize}
        \setlength{\itemsep}{1pt}
        \item Conmutativa: $a \cdot b = b \cdot a$
        \item Asociativa: $(a \cdot b) \cdot c = a \cdot (b \cdot c)$
        \item Neutro: $a \cdot 1 = a$ \quad$\cdot$\quad $a \cdot 0 = 0$
      \end{itemize}
    \end{mwextra}
  \end{columns}
  \begin{alertblock}{Regla de signos del producto}
    Positivo si el número de factores negativos es \textbf{par};
    \textbf{negativo} si es \textbf{impar}:
    $(-)\cdot(-) = (+)$, \; $(+)\cdot(-) = (-)$.
  \end{alertblock}
\end{frame}

\begin{frame}{Números Racionales y Decimales}
  \begin{block}{Conjunto de los Números Racionales}
    \[
      \mathbb{Q} = \left\{\frac{a}{b} \;\middle|\; a \in \mathbb{Z},\;
      b \in \mathbb{Z}^{*}\right\}
    \]
    $a \to$ \textbf{numerador}, \quad $b \to$ \textbf{denominador}.
    Todo número de $\mathbb{Z}$ es también racional:
    $\mathbb{N} \subset \mathbb{Z} \subset \mathbb{Q}$.
  \end{block}
  \begin{block}{Tipos de números decimales}
    \renewcommand{\arraystretch}{1.05}
    \begin{tabular}{@{} p{3.8cm} p{6.4cm} p{2.8cm} @{}}
      \toprule
      \textbf{Tipo} & \textbf{Descripción} & \textbf{Ejemplo} \\
      \midrule
      Limitados & Terminan & $0{,}5$; $1{,}25$ \\
      Periódicos puros & Ilimitado; el período empieza en la 1.ª
        cifra decimal & $0{,}\overline{3}$ \\
      Periódicos mixtos & Ilimitado; hay cifras antes del período &
        $0{,}1\overline{6}$ \\
      \bottomrule
    \end{tabular}
  \end{block}
\end{frame}

\begin{frame}{Ejercicios Propuestos — Conjuntos Numéricos \textnormal{(1/3)}}
  \begin{columns}[T]
    \column{0.48\textwidth}
    \begin{mwejercicio}[Ejercicio 1 \quad \facil]
      Clasifica cada número en $\mathbb{N}$, $\mathbb{Z}$ y $\mathbb{Q}$:
      \[
        0,\quad 5,\quad -12,\quad \frac{3}{4},\quad
        -\frac{7}{2},\quad 0{,}25
      \]
    \end{mwejercicio}
    \column{0.48\textwidth}
    \begin{mwejercicio}[Ejercicio 2 \quad \facil]
      Completa:
      \[
        \mathbb{Z}^{+} = \text{\underline{\hspace{1.2cm}}},
        \quad
        \mathbb{Z}^{-} = \text{\underline{\hspace{1.2cm}}}
      \]
      \[
        \mathbb{Z}^{*} = \text{\underline{\hspace{1.2cm}}}
      \]
      ¿Es cierto que $\mathbb{N} \subset \mathbb{Z} \subset \mathbb{Q}$?
      Justifica con ejemplos.
    \end{mwejercicio}
  \end{columns}
\end{frame}

\begin{frame}{Ejercicios Propuestos — Conjuntos Numéricos \textnormal{(2/3)}}
  \begin{columns}[T]
    \column{0.48\textwidth}
    \begin{mwejercicio}[Ejercicio 3 — Signos \quad \medio]
      Calcula aplicando la regla de signos:
      \begin{enumerate}
        \item $(-2)\cdot(-3)\cdot(-4)$
        \item $(-5)\cdot 6 \cdot (-2)$
        \item $(-1)\cdot(-1)\cdot(-1)\cdot(-1)$
      \end{enumerate}
    \end{mwejercicio}
    \column{0.48\textwidth}
    \begin{mwejercicio}[Ejercicio 4 — Propiedades \quad \medio]
      ¿Qué propiedad se está usando?
      \begin{enumerate}
        \item $3 + 7 = 7 + 3$
        \item $(2 + 3) + 4 = 2 + (3 + 4)$
        \item $5 + 0 = 5$
        \item $8 + (-8) = 0$
      \end{enumerate}
    \end{mwejercicio}
  \end{columns}
\end{frame}

\begin{frame}{Ejercicios Propuestos — Conjuntos Numéricos \textnormal{(3/3)}}
  \begin{mwejercicio}[Ejercicio 5 — Decimales \quad \medio]
    Clasifica cada número decimal como \textbf{limitado},
    \textbf{ilimitado periódico puro} o \textbf{ilimitado periódico
    mixto}:
    \[
      0{,}5 \qquad 0{,}333\dots \qquad 0{,}1666\dots
      \qquad 2{,}75 \qquad 4{,}2727\dots
    \]
    Escribe los periódicos usando la notación con raya (barra).
  \end{mwejercicio}
\end{frame}

% ════════════════════════════════════════════════════════════
\section{Potenciación en \texorpdfstring{$\mathbb{Q}$}{Q}}
% ════════════════════════════════════════════════════════════

\begin{frame}{Potenciación de una Fracción}
  \begin{block}{Definición}
    \[
      \left(\frac{a}{b}\right)^{n}
      = \frac{a^{n}}{b^{n}}
      = \underbrace{\frac{a}{b} \cdot \frac{a}{b} \cdots
        \frac{a}{b}}_{n \text{ veces}},
      \qquad n \in \mathbb{N},\; n \ge 1
    \]
  \end{block}
  \vspace{4pt}
  \begin{exampleblock}{Ejemplo}
    \[
      \left(\frac{2}{3}\right)^{3}
      = \frac{2^{3}}{3^{3}}
      = \frac{8}{27}
    \]
    Se eleva el numerador y el denominador al mismo exponente.
  \end{exampleblock}
\end{frame}

\begin{frame}{Propiedades de la Potenciación}
  \renewcommand{\arraystretch}{1.35}
  \begin{tabular}{@{} p{6.6cm} l @{}}
    \toprule
    \textbf{Propiedad} & \textbf{Expresión} \\
    \midrule
    Potencia cero &
      $\left(\frac{a}{b}\right)^{0} = 1$ \\
    Potencia uno &
      $\left(\frac{a}{b}\right)^{1} = \frac{a}{b}$ \\
    Producto de potencias (misma base) &
      $\left(\frac{a}{b}\right)^{n} \cdot
       \left(\frac{a}{b}\right)^{m}
       = \left(\frac{a}{b}\right)^{n+m}$ \\
    Cociente de potencias (misma base) &
      $\dfrac{\left(\frac{a}{b}\right)^{n}}
            {\left(\frac{a}{b}\right)^{m}}
       = \left(\frac{a}{b}\right)^{n-m}$ \\
    Potencia de una potencia &
      $\left[\left(\frac{a}{b}\right)^{n}\right]^{m}
       = \left(\frac{a}{b}\right)^{n \cdot m}$ \\
    Potencia del inverso &
      $\left(\frac{a}{b}\right)^{-n}
       = \left(\frac{b}{a}\right)^{n}$ \\
    \bottomrule
  \end{tabular}
\end{frame}

\begin{frame}{Ejercicios Propuestos — Potenciación \textnormal{(1/2)}}
  \begin{columns}[T]
    \column{0.48\textwidth}
    \begin{mwejercicio}[Ejercicio 1 \quad \facil]
      Calcula:
      \[
        \left(\frac{2}{5}\right)^{2},
        \qquad
        \left(-\frac{3}{4}\right)^{2},
        \qquad
        \left(\frac{1}{2}\right)^{4}
      \]
    \end{mwejercicio}
    \column{0.48\textwidth}
    \begin{mwejercicio}[Ejercicio 2 \quad \facil]
      Aplica potencia cero y potencia uno:
      \[
        \left(\frac{7}{9}\right)^{0},
        \qquad
        \left(-\frac{5}{6}\right)^{1},
        \qquad
        \left(\frac{2}{3}\right)^{0} +
        \left(\frac{2}{3}\right)^{1}
      \]
    \end{mwejercicio}
  \end{columns}
\end{frame}

\begin{frame}{Ejercicios Propuestos — Potenciación \textnormal{(2/2)}}
  \begin{columns}[T]
    \column{0.48\textwidth}
    \begin{mwejercicio}[Ejercicio 3 \quad \medio]
      Simplifica usando las propiedades:
      \[
        \left(\frac{2}{3}\right)^{3} \cdot
        \left(\frac{2}{3}\right)^{2},
        \qquad
        \frac{\left(\frac{2}{3}\right)^{5}}
             {\left(\frac{2}{3}\right)^{2}},
        \qquad
        \left[\left(\frac{3}{4}\right)^{2}\right]^{3}
      \]
    \end{mwejercicio}
    \column{0.48\textwidth}
    \begin{mwejercicio}[Ejercicio 4 \quad \medio]
      Escribe con exponente \textbf{positivo} y luego calcula:
      \[
        \left(\frac{2}{3}\right)^{-2},
        \qquad
        \left(\frac{5}{4}\right)^{-1},
        \qquad
        \left(\frac{1}{2}\right)^{-3}
      \]
    \end{mwejercicio}
  \end{columns}
\end{frame}

% ════════════════════════════════════════════════════════════
\section{Suma y Resta de Fracciones}
% ════════════════════════════════════════════════════════════

\begin{frame}{Fracciones con Igual Denominador}
  \begin{block}{Regla}
    \[
      \frac{a}{b} + \frac{c}{b} = \frac{a + c}{b}
    \]
    \textbf{Ejemplo:}\quad
    $\dfrac{3}{5} + \dfrac{1}{5} = \dfrac{4}{5}$
  \end{block}
  \vspace{4pt}
  \begin{mwextra}[Propiedades que se cumplen en $\mathbb{Q}$]
    \begin{enumerate}
      \item \textbf{Conmutativa:} $\frac{a}{b} + \frac{c}{d}
            = \frac{c}{d} + \frac{a}{b}$
      \item \textbf{Asociativa:} $\left(\frac{a}{b} +
            \frac{c}{d}\right) + \frac{e}{f}
            = \frac{a}{b} + \left(\frac{c}{d} + \frac{e}{f}\right)$
      \item \textbf{Elemento neutro:} $\frac{a}{b} + 0 = \frac{a}{b}$
      \item \textbf{Elemento inverso (opuesto):}
            $\frac{a}{b} + \left(-\frac{a}{b}\right) = 0$
    \end{enumerate}
  \end{mwextra}
\end{frame}

\begin{frame}{Denominadores Distintos: el M.C.M.}
  \begin{mwpasos}[Procedimiento]
    \begin{enumerate}
      \item Calcular el \textbf{m.c.m.} de los denominadores.
      \item Convertir cada fracción en una equivalente con ese
            denominador.
      \item Sumar (o restar) los numeradores y mantener el
            denominador.
      \item Simplificar el resultado si es posible.
    \end{enumerate}
  \end{mwpasos}
  \vspace{4pt}
  \begin{exampleblock}{Ejemplo de clase}
    \[
      \frac{2}{5} - \left(\frac{1}{3} + \frac{1}{2}\right)
      = \frac{12 - (10 + 15)}{30}
      = \frac{12 - 25}{30}
      = -\frac{13}{30}
    \]
    $\text{m.c.m.}\{5,\, 3,\, 2\} = 30$
  \end{exampleblock}
\end{frame}

\begin{frame}{Ejercicios Propuestos — Suma y Resta \textnormal{(1/3)}}
  \begin{columns}[T]
    \column{0.48\textwidth}
    \begin{mwejercicio}[Ejercicio 1 — Igual denominador \quad \facil]
      Calcula:
      \[
        \frac{3}{7} + \frac{2}{7},
        \qquad
        \frac{5}{9} - \frac{7}{9}
      \]
    \end{mwejercicio}
    \column{0.48\textwidth}
    \begin{mwejercicio}[Ejercicio 2 — Distinto denominador \quad \facil]
      Calcula usando el m.c.m.:
      \[
        \frac{1}{2} + \frac{1}{3},
        \qquad
        \frac{3}{4} - \frac{2}{5}
      \]
    \end{mwejercicio}
  \end{columns}
\end{frame}

\begin{frame}{Ejercicios Propuestos — Suma y Resta \textnormal{(2/3)}}
  \begin{columns}[T]
    \column{0.48\textwidth}
    \begin{mwejercicio}[Ejercicio 3 \quad \medio]
      Calcula y simplifica:
      \[
        \frac{7}{12} + \frac{5}{18} - \frac{1}{4}
      \]
      \textit{Pista:} $\text{m.c.m.}\{12,\, 18,\, 4\} = 36$.
    \end{mwejercicio}
    \column{0.48\textwidth}
    \begin{mwejercicio}[Ejercicio 4 \quad \medio]
      Resuelve respetando la jerarquía de operaciones:
      \[
        \frac{2}{5} - \left(\frac{1}{3} + \frac{1}{2}\right)
      \]
      \textit{Pista:} $\text{m.c.m.}\{5,\, 3,\, 2\} = 30$.
    \end{mwejercicio}
  \end{columns}
\end{frame}

\begin{frame}{Ejercicios Propuestos — Suma y Resta \textnormal{(3/3)}}
  \begin{mwejercicio}[Ejercicio 5 — Elementos neutro e inverso \quad \medio]
    \begin{enumerate}
      \item Escribe el \textbf{elemento inverso} de $\dfrac{3}{4}$ y
            verifica que la suma de ambos da $0$.
      \item ¿Qué número debes sumar a $-\dfrac{7}{10}$ para obtener
            $0$?
      \item Demuestra que $0$ es el elemento neutro usando
            $\dfrac{5}{8}$.
    \end{enumerate}
  \end{mwejercicio}
\end{frame}

% ════════════════════════════════════════════════════════════
\section{Producto y División de Fracciones}
% ════════════════════════════════════════════════════════════

\begin{frame}{Producto y División de Fracciones}
  \begin{columns}[T]
    \column{0.48\textwidth}
    \begin{block}{Producto}
      \[
        \frac{a}{b} \cdot \frac{c}{d}
        = \frac{a \cdot c}{b \cdot d}
      \]
      \textbf{Ejemplo:}
      \[
        \frac{3}{5} \cdot \left(-\frac{1}{9}\right)
        = \frac{-3}{45} = -\frac{1}{15}
      \]
    \end{block}
    \column{0.48\textwidth}
    \begin{block}{División}
      \[
        \frac{a}{b} \div \frac{c}{d}
        = \frac{a \cdot d}{b \cdot c}
      \]
      \textbf{Ejemplo:}
      \[
        \frac{4}{5} \div \frac{2}{15}
        = \frac{4 \cdot 15}{5 \cdot 2}
        = \frac{60}{10} = 6
      \]
    \end{block}
  \end{columns}
  \vspace{4pt}
  \begin{mwpasos}[Consejo]
    Multiplica los numeradores y los denominadores; al final,
    \textbf{simplifica} la fracción resultante hasta que sea
    irreducible.
  \end{mwpasos}
\end{frame}

\begin{frame}{Ejercicios Propuestos — Producto y División \textnormal{(1/2)}}
  \begin{columns}[T]
    \column{0.48\textwidth}
    \begin{mwejercicio}[Ejercicio 1 — Producto \quad \facil]
      Multiplica y simplifica:
      \[
        \frac{3}{5} \cdot \left(-\frac{1}{9}\right),
        \quad
        \frac{7}{10} \cdot \frac{5}{21},
        \quad
        \left(-\frac{2}{3}\right) \cdot \left(-\frac{9}{8}\right)
      \]
    \end{mwejercicio}
    \column{0.48\textwidth}
    \begin{mwejercicio}[Ejercicio 2 — División \quad \facil]
      Divide y simplifica:
      \[
        \frac{4}{5} \div \frac{2}{15},
        \quad
        \frac{2}{3} \div \frac{4}{9},
        \quad
        \frac{5}{6} \div \frac{10}{3}
      \]
    \end{mwejercicio}
  \end{columns}
\end{frame}

\begin{frame}{Ejercicios Propuestos — Producto y División \textnormal{(2/2)}}
  \begin{columns}[T]
    \column{0.48\textwidth}
    \begin{mwejercicio}[Ejercicio 3 — Mixtas \quad \medio]
      Resuelve respetando la jerarquía de operaciones:
      \[
        \left(\frac{3}{4} \div \frac{3}{8}\right) \cdot \frac{2}{5}
      \]
      \[
        \frac{1}{2} \cdot \frac{3}{4} + \frac{1}{3} \cdot \frac{9}{10}
      \]
    \end{mwejercicio}
    \column{0.48\textwidth}
    \begin{mwejercicio}[Ejercicio 4 — Problema \quad \medio]
      Para una receta se necesitan $\dfrac{3}{4}$ de taza de harina.
      \begin{enumerate}
        \item ¿Cuánta harina se necesita para $\dfrac{1}{2}$ receta?
        \item ¿Y para 3 recetas?
      \end{enumerate}
    \end{mwejercicio}
  \end{columns}
\end{frame}

% ════════════════════════════════════════════════════════════
\section{Fracciones Generatrices}
% ════════════════════════════════════════════════════════════

\begin{frame}{Decimales Limitados}
  \begin{block}{Definición}
    Para pasar un decimal \textbf{limitado} a fracción: coloca el
    número completo \textbf{sin la coma} en el numerador y la unidad
    seguida de \textbf{tantos ceros como cifras decimales} tenga en el
    denominador.
  \end{block}
  \vspace{4pt}
  \begin{exampleblock}{Ejemplos}
    \[
      0{,}145 = \frac{145}{1000} = \frac{29}{200}
      \qquad\qquad
      4{,}75 = \frac{475}{100} = \frac{19}{4}
    \]
  \end{exampleblock}
  \vspace{4pt}
  \begin{mwpasos}[Clave]
    La fracción generatriz \textbf{equivale exactamente} al decimal
    original: $29/200 = 0{,}145$.
  \end{mwpasos}
\end{frame}

\begin{frame}{Ilimitados Periódicos Puros}
  \begin{block}{Regla}
    El período empieza justo después de la coma.
    \begin{itemize}
      \item \textbf{Numerador:} el número sin la coma, \textbf{menos}
            la parte entera.
      \item \textbf{Denominador:} tantos \textbf{9} como cifras tenga
            el período.
    \end{itemize}
  \end{block}
  \vspace{4pt}
  \begin{exampleblock}{Ejemplos}
    \[
      13{,}\overline{9} = \frac{139 - 13}{9} = \frac{126}{9} = 14
      \qquad\qquad
      0{,}\overline{36} = \frac{36}{99} = \frac{4}{11}
    \]
  \end{exampleblock}
\end{frame}

\begin{frame}{Ilimitados Periódicos Mixtos}
  \begin{block}{Regla}
    Hay cifras decimales \textbf{antes} del período (parte
    anteperiódica).
    \begin{itemize}
      \item \textbf{Numerador:} el número sin la coma, \textbf{menos}
            las cifras no periódicas (parte entera + decimales
            anteperiódicos).
      \item \textbf{Denominador:} tantos \textbf{9} como cifras tenga
            el período y tantos \textbf{0} como cifras decimales
            anteperiódicas.
    \end{itemize}
  \end{block}
  \vspace{4pt}
  \begin{exampleblock}{Ejemplos}
    \[
      1{,}3\overline{5} = \frac{135 - 13}{90} = \frac{122}{90}
        = \frac{61}{45}
      \qquad
      31{,}0\overline{36} = \frac{31036 - 310}{990}
        = \frac{30726}{990} = \frac{1707}{55}
    \]
  \end{exampleblock}
\end{frame}

\begin{frame}{Ejercicios Propuestos — Fracciones Generatrices \textnormal{(1/2)}}
  \begin{columns}[T]
    \column{0.48\textwidth}
    \begin{mwejercicio}[Ejercicio 1 — Limitados \quad \facil]
      Pasa a fracción generatriz y simplifica:
      \[
        0{,}25 \qquad 0{,}145 \qquad 3{,}2
      \]
    \end{mwejercicio}
    \column{0.48\textwidth}
    \begin{mwejercicio}[Ejercicio 2 — Periódicos puros \quad \medio]
      Pasa a fracción generatriz y simplifica:
      \[
        0{,}\overline{36} \qquad 5{,}\overline{18}
        \qquad 13{,}\overline{9}
      \]
    \end{mwejercicio}
  \end{columns}
\end{frame}

\begin{frame}{Ejercicios Propuestos — Fracciones Generatrices \textnormal{(2/2)}}
  \begin{columns}[T]
    \column{0.48\textwidth}
    \begin{mwejercicio}[Ejercicio 3 — Periódicos mixtos \quad \dificil]
      Pasa a fracción generatriz y simplifica:
      \[
        1{,}3\overline{5} \qquad 2{,}1\overline{6}
        \qquad 31{,}0\overline{36}
      \]
    \end{mwejercicio}
    \column{0.48\textwidth}
    \begin{mwejercicio}[Ejercicio 4 — Ida y vuelta \quad \medio]
      ¿Qué decimal periódico representa $\dfrac{4}{11}$?
      Comprueba tu resultado convirtiendo el decimal de vuelta a
      fracción.
    \end{mwejercicio}
  \end{columns}
\end{frame}

% ════════════════════════════════════════════════════════════
\section{Fracción Irreducible}
% ════════════════════════════════════════════════════════════

\begin{frame}{¿Cuándo es Irreducible una Fracción?}
  \begin{block}{Definición}
    Una fracción es \textbf{irreducible} cuando su numerador y su
    denominador \textbf{no tienen divisores comunes} (distintos de 1).
    Si los tienen, es \textbf{reducible}: se simplifica dividiendo
    ambos por el común divisor.
  \end{block}
  \vspace{4pt}
  \begin{exampleblock}{Ejemplos de simplificación}
    \[
      \frac{33}{27} = \frac{33 \div 3}{27 \div 3} = \frac{11}{9}
      \qquad
      \frac{22}{72} = \frac{22 \div 2}{72 \div 2} = \frac{11}{36}
      \qquad
      \frac{42}{54} = \frac{14}{18} = \frac{7}{9}
    \]
    \[
      \frac{6}{8} = \frac{3}{4} = 0{,}75
    \]
  \end{exampleblock}
\end{frame}

\begin{frame}{De Fracción a Decimal}
  \begin{mwpasos}[Procedimiento]
    Divide el numerador entre el denominador. El resultado es un
    decimal \textbf{limitado} o \textbf{periódico}, según el tipo de
    fracción.
  \end{mwpasos}
  \vspace{4pt}
  \begin{exampleblock}{Verificación con fracciones generatrices}
    \[
      \frac{4}{3} = 1{,}\overline{3}
      \qquad
      \frac{5}{6} = 0{,}8\overline{3}
      \qquad
      \frac{1}{10} = 0{,}1
    \]
  \end{exampleblock}
  \vspace{4pt}
  \begin{mwextra}[Comprobación inversa]
    \[
      1{,}\overline{3} = \frac{13 - 1}{9} = \frac{12}{9} = \frac{4}{3}
      \quad \checkmark
    \]
  \end{mwextra}
\end{frame}

\begin{frame}{Ejercicios Propuestos — Fracción Irreducible \textnormal{(1/2)}}
  \begin{columns}[T]
    \column{0.48\textwidth}
    \begin{mwejercicio}[Ejercicio 1 \quad \facil]
      Simplifica hasta llegar a la fracción irreducible:
      \[
        \frac{33}{27}, \quad \frac{22}{72}, \quad
        \frac{42}{54}, \quad \frac{45}{60}
      \]
    \end{mwejercicio}
    \column{0.48\textwidth}
    \begin{mwejercicio}[Ejercicio 2 \quad \facil]
      ¿Cuáles de estas fracciones ya son irreducibles?
      \[
        \frac{4}{5}, \quad \frac{7}{9}, \quad
        \frac{13}{20}, \quad \frac{6}{8}
      \]
    \end{mwejercicio}
  \end{columns}
\end{frame}

\begin{frame}{Ejercicios Propuestos — Fracción Irreducible \textnormal{(2/2)}}
  \begin{columns}[T]
    \column{0.48\textwidth}
    \begin{mwejercicio}[Ejercicio 3 \quad \medio]
      Escribe cada fracción como número decimal:
      \[
        \frac{4}{3}, \quad \frac{5}{6}, \quad
        \frac{1}{10}, \quad \frac{3}{8}
      \]
    \end{mwejercicio}
    \column{0.48\textwidth}
    \begin{mwejercicio}[Ejercicio 4 \quad \medio]
      Verifica, usando fracciones generatrices, que se cumple:
      \[
        \frac{4}{3} = 1{,}\overline{3},
        \quad
        \frac{5}{6} = 0{,}8\overline{3},
        \quad
        \frac{1}{10} = 0{,}1
      \]
    \end{mwejercicio}
  \end{columns}
\end{frame}

% ════════════════════════════════════════════════════════════
% ============================================================
\section{Números Reales y Aproximaciones}
% ============================================================

\begin{frame}{El Conjunto de los Números Reales}
  \begin{block}{Definición}
    \[
      \mathbb{R} = \mathbb{Q} \cup \mathbb{I}
      \qquad\qquad
      \mathbb{N} \subset \mathbb{Z} \subset \mathbb{Q} \subset \mathbb{R}
    \]
  \end{block}
  \vspace{2pt}
  \begin{columns}[T]
    \column{0.48\textwidth}
    \begin{exampleblock}{Racionales ($\mathbb{Q}$)}
      Se escriben como $\dfrac{a}{b}$: decimales \textbf{limitados} o
      \textbf{periódicos}.
      \[
        \frac{3}{4} = 0{,}75, \quad
        \frac{1}{3} = 0{,}\overline{3}
      \]
    \end{exampleblock}
    \column{0.48\textwidth}
    \begin{mwextra}[Irracionales ($\mathbb{I}$)]
      Decimales \textbf{infinitos no periódicos}: \textbf{no} se
      pueden escribir como fracción.
      \[
        \sqrt{2} = 1{,}4142\dots, \quad
        \pi = 3{,}1415\dots
      \]
    \end{mwextra}
  \end{columns}
  \vspace{2pt}
  \begin{mwpasos}[Recta real]
    Cada punto de la recta numérica representa un número real y
    viceversa: $\mathbb{R}$ \textbf{llena} la recta.
  \end{mwpasos}
\end{frame}

\begin{frame}{Aproximación por Exceso y por Defecto}
  \begin{columns}[T]
    \column{0.48\textwidth}
    \begin{block}{Por exceso}
      El valor de precisión \textbf{aumenta una unidad}: se corta el
      número y se suma $1$ a la última cifra conservada.
      \vspace{2pt}

      \textbf{Ejemplo:} $0{,}4325$ a la milésima
      \[
        0{,}4325 \approx 0{,}433
      \]
      La cifra siguiente ($5$) obliga a subir.
    \end{block}
    \column{0.48\textwidth}
    \begin{mwextra}[Por defecto]
      Se deja \textbf{igual} el valor de precisión y se \textbf{corta}
      el número (truncamiento).
      \vspace{2pt}

      \textbf{Ejemplo:} $1{,}4324$ a la centésima
      \[
        1{,}4324 \approx 1{,}43
      \]
      Se conservan $2$ cifras decimales.
    \end{mwextra}
  \end{columns}
  \begin{alertblock}{El mismo número, dos resultados}
    $0{,}4325$ a la milésima: por \textbf{defecto} $0{,}432$
    \quad$\cdot$\quad por \textbf{exceso} $0{,}433$.
  \end{alertblock}
\end{frame}

\begin{frame}{Ejercicios Propuestos — Aproximaciones \textnormal{(1/2)}}
  \begin{columns}[T]
    \column{0.48\textwidth}
    \begin{mwejercicio}[Ejercicio 1 — Por exceso \quad \facil]
      Aproxima a la precisión indicada \textbf{por exceso}:
      \begin{enumerate}
        \item $0{,}4325$ a la milésima
        \item $2{,}7381$ a la centésima
        \item $\pi = 3{,}14159\dots$ a la diezmilésima
      \end{enumerate}
    \end{mwejercicio}
    \column{0.48\textwidth}
    \begin{mwejercicio}[Ejercicio 2 — Por defecto \quad \facil]
      Aproxima a la precisión indicada \textbf{por defecto}:
      \begin{enumerate}
        \item $1{,}4324$ a la centésima
        \item $7{,}9998$ a la décima
        \item $\sqrt{2} = 1{,}41421\dots$ a la milésima
      \end{enumerate}
    \end{mwejercicio}
  \end{columns}
\end{frame}

\begin{frame}{Ejercicios Propuestos — Aproximaciones \textnormal{(2/2)}}
  \begin{mwejercicio}[Ejercicio 3 — Tabla comparativa \quad \medio]
    Completa la tabla con la aproximación de cada número en la
    precisión indicada:
    \vspace{4pt}

    \renewcommand{\arraystretch}{1.3}
    \begin{tabular}{@{} l c c c @{}}
      \toprule
      \textbf{Número} & \textbf{Precisión} & \textbf{Por defecto}
        & \textbf{Por exceso} \\
      \midrule
      $0{,}4325$ & milésima & & \\
      $1{,}4324$ & centésima & & \\
      $3{,}14159\dots$ & diezmilésima & & \\
      $9{,}8765$ & décima & & \\
      \bottomrule
    \end{tabular}
  \end{mwejercicio}
\end{frame}

\begin{frame}{Soluciones — Aproximaciones \textnormal{(1/2)}}
  \mwtopanchoroff
  \begin{columns}[T]
    \column{0.48\textwidth}
    \begin{exampleblock}{Ejercicio 1 — Por exceso}
      \begin{itemize}
        \item $0{,}4325 \approx 0{,}433$
        \item $2{,}7381 \approx 2{,}74$
        \item $\pi \approx 3{,}1416$
      \end{itemize}
      \small (se corta y se sube una unidad en la última cifra)
    \end{exampleblock}
    \column{0.48\textwidth}
    \begin{mwextra}[Ejercicio 2 — Por defecto]
      \begin{itemize}
        \item $1{,}4324 \approx 1{,}43$
        \item $7{,}9998 \approx 7{,}9$
        \item $\sqrt{2} \approx 1{,}414$
      \end{itemize}
      \small (solo se corta: las cifras restantes se descartan)
    \end{mwextra}
  \end{columns}
  \mwtopanchoron
\end{frame}

\begin{frame}{Soluciones — Aproximaciones \textnormal{(2/2)}}
  \mwtopanchoroff
  \begin{block}{Ejercicio 3 — Tabla completa}
    \renewcommand{\arraystretch}{1.25}
    \begin{tabular}{@{} l c c c @{}}
      \toprule
      \textbf{Número} & \textbf{Precisión} & \textbf{Por defecto}
        & \textbf{Por exceso} \\
      \midrule
      $0{,}4325$ & milésima & $0{,}432$ & $0{,}433$ \\
      $1{,}4324$ & centésima & $1{,}43$ & $1{,}44$ \\
      $3{,}14159\dots$ & diezmilésima & $3{,}1415$ & $3{,}1416$ \\
      $9{,}8765$ & décima & $9{,}8$ & $9{,}9$ \\
      \bottomrule
    \end{tabular}
  \end{block}
  \vspace{2pt}
  \begin{mwpasos}[Cómo revisar tu tabla]
    Cuenta las cifras decimales de la precisión: debe haber exactamente
    esas cifras. La diferencia entre ambas columnas es \textbf{una
    unidad} en la última posición.
  \end{mwpasos}
  \mwtopanchoron
\end{frame}

% ============================================================
\section{Práctica: Operaciones Combinadas}
% ============================================================

\begin{frame}{Operaciones Combinadas: Plan de Ataque}
  \begin{mwpasos}[Procedimiento]
    \begin{enumerate}
      \item Convierte \textbf{cada decimal} a fracción generatriz
            (limitado, periódico puro o mixto).
      \item Simplifica cada fracción antes de operar.
      \item Respeta la \textbf{jerarquía}: paréntesis $\to$
            potencias $\to$ $\times$ y $\div$ $\to$ $+$ y $-$.
      \item Lleva el resultado a \textbf{fracción irreducible} y, si se
            pide, a decimal.
    \end{enumerate}
  \end{mwpasos}
  \begin{alertblock}{Cuidado con la división de decimales}
    $6{,}2 \div 3{,}1$: convierte primero a fracciones
    $\dfrac{62}{10} \div \dfrac{31}{10} = \dfrac{62}{31} = 2$.
  \end{alertblock}
  \vspace{2pt}
  \begin{mwextra}[Generatrices que necesitarás]
    $0{,}2 = \dfrac{2}{10}$ \quad$\cdot$\quad
    $0{,}6 = \dfrac{6}{10}$ \quad$\cdot$\quad
    $0{,}\overline{3} = \dfrac{3}{9}$ \quad$\cdot$\quad
    $0{,}\overline{5} = \dfrac{5}{9}$ \quad$\cdot$\quad
    $3{,}\overline{9} = \dfrac{39 - 3}{9}$
  \end{mwextra}
\end{frame}

\begin{frame}{Ejercicio Resuelto — Del Cuaderno}
  \mwtopanchoroff
  \begin{mwpasos}[Calcula el valor exacto]
    \textbf{Enunciado: } $1 - 6{,}2 \div 3{,}1 + (0{,}6)^{2}$.
    \begin{align*}
      (0{,}6)^{2} &= 0{,}36 = \frac{36}{100} = \frac{9}{25} \\
      6{,}2 \div 3{,}1
        &= \frac{62}{10} \div \frac{31}{10}
         = \frac{62}{31} = 2 \\
      1 - 6{,}2 \div 3{,}1 + (0{,}6)^{2}
        &= 1 - 2 + 0{,}36 = -1 + 0{,}36 \\
        &= -0{,}64 = -\frac{64}{100} = -\frac{16}{25}
    \end{align*}
  \end{mwpasos}
  \mwtopanchoron
\end{frame}

\begin{frame}{Ejercicios Propuestos — Valor Exacto en $\mathbb{Q}$}
  \begin{mwejercicio}[Del cuaderno \quad \dificil]
    Convierte cada decimal en fracción generatriz y calcula el
    \textbf{valor exacto}:
    \begin{enumerate}
      \item[a)] $1 - 6{,}2 \div 3{,}1 + (0{,}6)^{2}$
      \item[b)] $\dfrac{0{,}\overline{3} - 0{,}\overline{5} - 0{,}1}
                       {3{,}5 - 2}$
      \item[c)] $\dfrac{6{,}5 + 2{,}7 + 3{,}\overline{9}}{4{,}5}$
    \end{enumerate}
  \end{mwejercicio}
\end{frame}

\begin{frame}{Solución Explicada — Ejercicio b)}
  \mwtopanchoroff
  \begin{block}{b) \; $\dfrac{0{,}\overline{3} - 0{,}\overline{5} - 0{,}1}
                                  {3{,}5 - 2}$}
    \textbf{Paso 1 — generatrices:}
    $0{,}\overline{3} = \frac{3}{9} = \frac{1}{3}$,\quad
    $0{,}\overline{5} = \frac{5}{9}$,\quad
    $0{,}1 = \frac{1}{10}$,\quad
    $3{,}5 - 2 = \frac{3}{2}$.
    \[
      \text{Paso 2 — numerador (m.c.m.\ } 90\text{):}\quad
      \frac{1}{3} - \frac{5}{9} - \frac{1}{10}
      = \frac{30 - 50 - 9}{90}
      = -\frac{29}{90}
    \]
    \[
      \text{Paso 3 — cociente:}\quad
      \frac{-29/90}{3/2}
      = -\frac{29}{90} \cdot \frac{2}{3}
      = -\frac{58}{270}
      = -\frac{29}{135}
    \]
  \end{block}
  \mwtopanchoron
\end{frame}

\begin{frame}{Solución Explicada — Ejercicio c)}
  \mwtopanchoroff
  \begin{block}{c) \; $\dfrac{6{,}5 + 2{,}7 + 3{,}\overline{9}}{4{,}5}$}
    \textbf{Paso 1 — el periódico puro:}
    $3{,}\overline{9} = \dfrac{39 - 3}{9} = 4$,
    \qquad $4{,}5 = \dfrac{9}{2}$.
    \[
      \text{Paso 2:}\quad 6{,}5 + 2{,}7 + 4 = 13{,}2 = \frac{132}{10}
      = \frac{66}{5}
    \]
    \[
      \text{Paso 3:}\quad \frac{66/5}{9/2}
      = \frac{66}{5} \cdot \frac{2}{9}
      = \frac{132}{45}
      = \frac{44}{15} \approx 2{,}9\overline{3}
    \]
    \small Comprobación: $13{,}2 \div 4{,}5 = 2{,}9333\dots
    = 2{,}9\overline{3}$ \;\checkmark
  \end{block}
  \mwtopanchoron
\end{frame}

\begin{frame}{Ejercicios Propuestos — Más Práctica \textnormal{(1/2)}}
  \begin{columns}[T]
    \column{0.48\textwidth}
    \begin{mwejercicio}[Ejercicio 1 \quad \medio]
      Calcula el valor exacto:
      \[
        1{,}5 + 0{,}\overline{3} \cdot 3
      \]
    \end{mwejercicio}
    \column{0.48\textwidth}
    \begin{mwejercicio}[Ejercicio 2 \quad \dificil]
      Calcula el valor exacto:
      \[
        \frac{0{,}4 - 0{,}\overline{1}}{1{,}\overline{2}}
      \]
    \end{mwejercicio}
  \end{columns}
\end{frame}

\begin{frame}{Ejercicios Propuestos — Más Práctica \textnormal{(2/2)}}
  \begin{columns}[T]
    \column{0.48\textwidth}
    \begin{mwejercicio}[Ejercicio 3 \quad \medio]
      Calcula el valor exacto:
      \[
        (2 - 1{,}\overline{3}) \div 0{,}5
      \]
    \end{mwejercicio}
    \column{0.48\textwidth}
    \begin{mwejercicio}[Ejercicio 4 \quad \medio]
      Escribe $\dfrac{1}{7}$ como decimal y aproxímalo a la centésima
      \textbf{por defecto} y \textbf{por exceso}. ¿Cuál está más cerca
      del valor real?
    \end{mwejercicio}
  \end{columns}
\end{frame}

\begin{frame}{Soluciones — Más Práctica \textnormal{(1/2)}}
  \mwtopanchoroff
  \begin{columns}[T]
    \column{0.48\textwidth}
    \begin{exampleblock}{Ejercicio 1}
      \[
        1{,}5 = \frac{3}{2}, \qquad
        0{,}\overline{3} \cdot 3 = \frac{1}{3} \cdot 3 = 1
      \]
      \[
        \frac{3}{2} + 1 = \frac{5}{2} = 2{,}5
      \]
    \end{exampleblock}
    \column{0.48\textwidth}
    \begin{mwextra}[Ejercicio 2]
      \[
        0{,}4 = \frac{2}{5}, \quad
        0{,}\overline{1} = \frac{1}{9}, \quad
        1{,}\overline{2} = \frac{11}{9}
      \]
      \[
        \frac{\frac{2}{5} - \frac{1}{9}}{\frac{11}{9}}
        = \frac{13}{45} \cdot \frac{9}{11}
        = \frac{13}{55}
      \]
    \end{mwextra}
  \end{columns}
  \mwtopanchoron
\end{frame}

\begin{frame}{Soluciones — Más Práctica \textnormal{(2/2)}}
  \mwtopanchoroff
  \begin{columns}[T]
    \column{0.48\textwidth}
    \begin{exampleblock}{Ejercicio 3}
      \[
        1{,}\overline{3} = \frac{12 - 1}{9} = \frac{4}{3},
        \qquad
        0{,}5 = \frac{1}{2}
      \]
      \[
        \frac{2 - \frac{4}{3}}{\frac{1}{2}}
        = \frac{2}{3} \cdot 2 = \frac{4}{3} = 1{,}\overline{3}
      \]
    \end{exampleblock}
    \column{0.48\textwidth}
    \begin{mwextra}[Ejercicio 4]
      \[
        \frac{1}{7} = 0{,}142857\dots
      \]
      \[
        \text{defecto: } 0{,}14 \qquad
        \text{exceso: } 0{,}15
      \]
      \small Está más cerca $0{,}14$: la diferencia es $0{,}0028\dots$
      frente a $0{,}0071\dots$
    \end{mwextra}
  \end{columns}
  \mwtopanchoron
\end{frame}

% ============================================================
\section{Ejercicios de Repaso General}
% ════════════════════════════════════════════════════════════

\begin{frame}{Cuadro Resumen — Jerarquía de Conjuntos}
  \[
    \mathbb{N} = \{0, 1, 2, \dots\}
    \;\subset\;
    \mathbb{Z} = \{\dots, -2, -1, 0, 1, \dots\}
    \;\subset\;
    \mathbb{Q} = \left\{\frac{a}{b} : a \in \mathbb{Z},\,
    b \in \mathbb{Z}^{*}\right\}
  \]
  \vspace{4pt}
  \begin{columns}[T]
    \column{0.32\textwidth}
    \begin{block}{Naturales}
      Enteros, discretos e infinitos.\\
      Neutro: $a + 0 = a$
    \end{block}
    \column{0.32\textwidth}
    \begin{mwextra}[Enteros]
      Incluyen los negativos.\\
      Inverso: $a + (-a) = 0$
    \end{mwextra}
    \column{0.32\textwidth}
    \begin{exampleblock}{Racionales}
      Fracciones $\dfrac{a}{b}$ y todos los decimales.
    \end{exampleblock}
  \end{columns}
  \vspace{2pt}
  \begin{center}\small
    Además $\mathbb{Q} \subset \mathbb{R} = \mathbb{Q} \cup \mathbb{I}$,
    donde $\mathbb{I}$ son los irracionales
    ($\sqrt{2}$, $\pi$, $e$): decimales infinitos no periódicos.
  \end{center}
\end{frame}

\begin{frame}{Repaso General \textnormal{(1/5)} — Conjuntos y Signos}
  \begin{mwejercicio}[Repaso 1 \quad \medio]
    \begin{enumerate}
      \item Clasifica cada número en $\mathbb{N}$, $\mathbb{Z}$ y
            $\mathbb{Q}$: $0$;\; $-7$;\; $\dfrac{22}{7}$;\;
            $0{,}25$;\; $15$.
      \item Calcula aplicando la regla de signos:
            $(-3)\cdot(-4)\cdot(-2)$ \; y \;
            $(-5)\cdot 6 \cdot (-1)$.
      \item Completa: $\mathbb{Z}^{+} = \_\_\_$,\;
            $\mathbb{Z}^{-} = \_\_\_$,\; $\mathbb{Z}^{*} = \_\_\_$.
    \end{enumerate}
  \end{mwejercicio}
\end{frame}

\begin{frame}{Repaso General \textnormal{(2/5)} — Potenciación}
  \begin{mwejercicio}[Repaso 2 \quad \medio]
    \begin{enumerate}
      \item Calcula: $\left(\dfrac{2}{5}\right)^{2}$,\;
            $\left(-\dfrac{3}{4}\right)^{3}$,\;
            $\left(\dfrac{7}{8}\right)^{0}$.
      \item Simplifica: $\left(\dfrac{2}{3}\right)^{4} \cdot
            \left(\dfrac{2}{3}\right)^{2}$,\;
            $\dfrac{\left(\frac{3}{5}\right)^{7}}
                   {\left(\frac{3}{5}\right)^{4}}$,\;
            $\left[\left(\dfrac{1}{2}\right)^{2}\right]^{3}$.
      \item Escribe con exponente positivo: $\left(\dfrac{3}{4}\right)^{-2}$
            y calcula su valor.
    \end{enumerate}
  \end{mwejercicio}
\end{frame}

\begin{frame}{Repaso General \textnormal{(3/5)} — Operaciones con Fracciones}
  \begin{mwejercicio}[Repaso 3 \quad \medio]
    Realiza las operaciones y simplifica:
    \[
      \text{a)}\;\; \frac{5}{12} + \frac{7}{18} - \frac{1}{4}
      \qquad
      \text{b)}\;\; \frac{3}{5} \cdot \left(-\frac{1}{9}\right)
    \]
    \[
      \text{c)}\;\; \frac{4}{5} \div \frac{2}{15}
      \qquad
      \text{d)}\;\; \frac{2}{5} - \left(\frac{1}{3} + \frac{1}{2}\right)
    \]
  \end{mwejercicio}
\end{frame}

\begin{frame}{Repaso General \textnormal{(4/5)} — Generatrices e Irreducibles}
  \begin{mwejercicio}[Repaso 4 \quad \dificil]
    \begin{enumerate}
      \item Pasa a fracción generatriz y simplifica: $0{,}145$;\;
            $1{,}3\overline{5}$;\; $31{,}0\overline{36}$.
      \item Escribe como número decimal: $\dfrac{4}{3}$;\;
            $\dfrac{5}{6}$;\; $\dfrac{1}{10}$.
      \item Simplifica hasta la fracción irreducible:
            $\dfrac{33}{27}$,\; $\dfrac{22}{72}$,\; $\dfrac{42}{54}$.
    \end{enumerate}
  \end{mwejercicio}
\end{frame}

\begin{frame}{Repaso General \textnormal{(5/5)} — Problemas}
  \begin{columns}[T]
    \column{0.48\textwidth}
    \begin{mwejercicio}[Repaso 5 — Galletas \quad \medio]
      Un paquete tiene 24 galletas. Ana come $\dfrac{1}{4}$ del
      paquete y Luis come $\dfrac{1}{6}$.
      \begin{enumerate}
        \item ¿Qué fracción del paquete comieron juntos?
        \item ¿Cuántas galletas quedan?
      \end{enumerate}
    \end{mwejercicio}
    \column{0.48\textwidth}
    \begin{mwejercicio}[Repaso 6 — Oferta \quad \medio]
      El precio de un libro baja $\dfrac{2}{9}$ de su valor. Si el
      libro costaba $\$9.000$:
      \begin{enumerate}
        \item ¿Cuánto baja el precio?
        \item ¿Cuál es el precio final?
      \end{enumerate}
    \end{mwejercicio}
  \end{columns}
\end{frame}

\end{document}
